\documentclass[10pt]{article}

\usepackage[T1]{fontenc}
\usepackage[utf8]{inputenc}
\usepackage{lmodern}
\usepackage[margin=0.85in]{geometry}
\usepackage{amsmath,amssymb}
\usepackage{array}
\usepackage{booktabs}
\usepackage{listings}
\usepackage{xcolor}
\usepackage{titlesec}
\usepackage[colorlinks=true,linkcolor=black,citecolor=black,urlcolor=blue]{hyperref}

\titlespacing*{\section}{0pt}{1.0ex plus .2ex minus .2ex}{0.6ex plus .1ex}
\titleformat{\section}{\normalfont\large\bfseries}{\thesection}{0.6em}{}
\setlength{\abovedisplayskip}{5pt}\setlength{\belowdisplayskip}{5pt}
\setlength{\abovedisplayshortskip}{3pt}\setlength{\belowdisplayshortskip}{3pt}
\setlength{\emergencystretch}{1em}

\newcommand{\ty}[1]{\mathtt{#1}}     % a type name, in monospace
\newcommand{\arr}{\to}               % "give ..., get ..."
\newcommand{\opt}{\mathbin{?}}       % option

\lstdefinestyle{eo}{%
  basicstyle=\ttfamily\small,
  columns=fullflexible,
  keepspaces=true,
  showstringspaces=false,
  breaklines=false,
  frame=none,
  xleftmargin=1.4em,
  aboveskip=0.3em,
  belowskip=0.3em}
\lstset{style=eo}

\title{\vspace{-1.5em}Toward Type Inference for EO\\[-0.2em]
  \large a working note: how decoration and $\bot$ already tell us the types}
\author{Objectionary\\\small\url{https://github.com/objectionary/eo}}
\date{July 2026}

\begin{document}
\maketitle
\vspace{-1.4em}

\begin{abstract}
\noindent
EO has no written-down types, but it wants a tool that catches type mistakes
before a program runs. This note explains, in plain terms, what ``type'' should
mean in EO, why the obvious idea (compare the \emph{name} of the returned
object) breaks the moment we use decoration, and what to use instead: a notion
of type based on an object's \emph{shape}. We then do three concrete things:
write down the grammar of these shapes, give the shape (the ``signature'') of
each built-in object, and describe --- as plain steps --- a checker that works
out the shape of any expression and reports an error when the shapes do not
fit. A worked example shows the checker accepting a correct program and
rejecting two broken ones. None of this is new type theory; the pieces are all
known and are pointed to in footnotes for the curious. The only EO-specific
parts are how we treat decoration and the failure value~$\bot$.
\end{abstract}

\section*{What you need to know (the whole vocabulary)}

Everything below uses just these words and this notation. If a footnote names a
paper, it is optional background, not required reading.

\begin{description}\setlength{\itemsep}{1pt}
\item[Type inference] the tool figures out the ``kind'' of every value by
  itself, because EO writes none down.\footnote{Figuring out types from
  unannotated code is \emph{type inference}; the foundational algorithm is
  L.~Damas \& R.~Milner, \emph{Principal type-schemes for functional programs},
  POPL 1982.}
\item[Nominal vs.\ structural] \emph{nominal}: two objects have the same type
  only if they share a \emph{name}. \emph{Structural}: they match if they have
  the same \emph{shape} --- the same attributes. EO needs structural.\footnote{%
  Treating a type as a shape (a set of attributes) rather than a name is
  standard; a classic reference is M.~Abadi \& L.~Cardelli, \emph{A Theory of
  Objects}, Springer 1996.}
\item[Usable-as (subtyping)] ``an X can go wherever a Y is wanted.'' If
  \texttt{my-num} has every attribute \texttt{number} has, then \texttt{my-num}
  is usable as a \texttt{number}.
\item[Notation] $X \arr Y$ means ``give it an $X$, get a $Y$'';\quad
  $A \mid B$ means ``an $A$ or a $B$'' (as in TypeScript);\quad
  $T\opt$ means ``a $T$, or $\bot$'' (the failure value);\quad
  $\{\,a:\dots,\ b:\dots\,\}$ is an object with attributes $a,b,\dots$;\quad
  ``for any $A$'' before a type means it works for every type (the tool picks
  which).
\end{description}

\section{The problem: EO names objects, but we care about their shape}

Every built-in EO object (an ``atom'') already declares the \emph{name} of what
it returns, after a \verb|/|. These are real lines from \texttt{number.eo}:
\begin{lstlisting}
[as-bytes] > number
  [x] > plus  /number      # returns a number
  [x] > times /number
  [x] > gt    /bool        # returns a bool
\end{lstlisting}
A name like \texttt{number} or \texttt{bool} is a \emph{nominal} tag. It works
until we decorate. Decoration is EO's ``\texttt{@}'' attribute: an object with
an \texttt{@} behaves like whatever its \texttt{@} points to, because a call it
does not answer itself is passed down to \texttt{@}. Consider:
\begin{lstlisting}
[] > my-num       # its name is `my-num'; it behaves like a number
  5 > @
\end{lstlisting}
\verb|my-num| is named \verb|my-num|, but calls like \verb|plus| or \verb|gt|
fall through \verb|@| to the \verb|5| inside, so it \emph{behaves} as a number.
A check that compares the name \verb|my-num| against the expected \verb|number|
wrongly says ``no.''

You might try to save the name-based idea by remembering the whole decoration
\emph{chain} (\verb|my-num| then \verb|number|) and checking whether
\verb|number| is in it. That survives the case above but not this one:
\begin{lstlisting}
[] > my-num       # @ is `seq *', whose name is `seq' -- `number' is nowhere
  seq * > @
    "Hello"
    5
\end{lstlisting}
Here \verb|my-num| still behaves like a number, because \verb|seq| runs its
steps and hands back the last one (\verb|5|). But knowing that needs the
\emph{meaning} of \verb|seq|, not its name. So the two examples mark a line: the
first is answerable by names, the second only by behaviour. Any name-based
scheme handles the first and fails the second --- so we have to describe
behaviour, i.e.\ \emph{shape}. This is also why EO was right to switch off its
old runtime check that compared names: the real check needs the object's shape,
and finding a shape at runtime means running the program. Shape is better
worked out ahead of time, by inference.

\section{What ``type'' should mean: a shape, and decoration means ``usable-as''}

Let a type describe what an object \emph{does}: which attributes it answers, and
with what. \verb|number|'s type is the shape
$\{\,\ty{plus}:\ty{number}\arr\ty{number},\ \ty{gt}:\ty{number}\arr\ty{bool},\ \dots\}$.
The one EO-specific rule is about \texttt{@}:

\begin{quote}
\textbf{Decoration rule.} If an object has \texttt{@} pointing to something of
shape $S$, then that object answers everything $S$ answers (directly, or by
falling through \texttt{@}). So it is \emph{usable as} an $S$.
\end{quote}

That is the whole trick: \verb|my-num| (which decorates \verb|5|) is usable as a
\verb|number| because its \texttt{@} is a number. The name \verb|number| does
not disappear --- it just becomes a convenient \emph{alias} for that shape, so
error messages can say ``number'' instead of printing the full attribute list.

\section{The type grammar (the shapes a type can take)}\label{sec:grammar}

A type is one of the following. Each line is a shape, with a tiny example.

\begin{description}\setlength{\itemsep}{2pt}
\item[An unknown, written $A$, $B$, \dots] a placeholder the tool has not filled
  in yet. A fresh unknown is created for each free attribute of a formation and
  gets pinned down as the tool learns more.
\item[An object shape $\{\,\ell_1:\tau_1,\ \dots,\ \ell_n:\tau_n\,\}$] a list of
  attribute names with their types. One attribute may be the special
  \texttt{@}. A trailing ``$\dots$'' means ``and possibly more attributes'' ---
  needed so a shape can stand for ``at least these.''\footnote{The ``at least
  these'' part is \emph{row polymorphism}: M.~Wand, LICS 1987; D.~R\'emy, 1994.}
\item[A tuple, $\ty{tuple}\ A$ or $\langle A, B, C\rangle$] an ordered
  collection; either ``a tuple of $A$'s'' or a fixed-length one with a type per
  position.
\item[Failure, $\bot$, and the option $T\opt$] $\bot$ is the value a terminated
  computation produces; $T\opt$ means ``a $T$, or $\bot$.''
\item[Either/or, $A \mid B$] a value that is an $A$ or a $B$ (for example, a
  branch that returns one of two things).
\item[A self-referring shape, ``loop $A.\ \dots$''] for objects that mention
  themselves, like a list whose tail is another list of the same
  kind.\footnote{\emph{Recursive types}; R.~Amadio \& L.~Cardelli, 1993.}
\end{description}

\emph{For the record (skippable), the same grammar in symbols:}
\[
\tau \;::=\; A \;\mid\; \{\ell_1{:}\tau_1,\dots,\ell_n{:}\tau_n\} \;\mid\;
\ty{tuple}\,\tau \;\mid\; \langle\tau,\dots,\tau\rangle \;\mid\; \bot \;\mid\;
\tau\opt \;\mid\; \tau \mid \tau \;\mid\; \mu A.\,\tau
\]

\section{Failure as a type: $\bot$ and \texttt{recovered}}\label{sec:fail}

EO recently made every unrecoverable failure produce one special value,
$\bot$, and the \emph{only} way to catch it is the object \verb|recovered|
(\texttt{recovered.eo}):
\begin{lstlisting}
[value alternative] > recovered /bytes      # value, unless it is a bottom,
recovered (2.as-i64.div 0.as-i64) 7         # in which case: alternative  ==> 7
\end{lstlisting}
In types, this is simple and powerful. An operation that might fail has type
$T\opt$ (``a $T$, or $\bot$''). Calling a normal attribute on a $T\opt$ is an
error --- you must clear the $\bot$ first. \verb|recovered| is what clears it:
give it a maybe-failed value and a fallback, and it hands back a definite value.
Its type is
\[
\ty{recovered} :\ \text{for any } A, B:\quad A\opt \arr B \arr (A \mid B),
\]
read as: ``take an $A$-or-$\bot$ and a $B$; give back an $A$ or a $B$'' (the
$\bot$ is gone). This is the same shape as ``get the value out of a box, or use
a default'' in other languages.\footnote{This need to combine ``a value or
nothing'' with ``an $A$ or a $B$'' in one inferred type is exactly what
\emph{algebraic subtyping} was built to do well: S.~Dolan, PhD thesis,
Cambridge 2017; and a very readable version, L.~Parreaux, \emph{The Simple
Essence of Algebraic Subtyping}, ICFP 2020 (about 500 lines of code).}

\paragraph{Two doors, and safe chaining with \texttt{?.}} Since a normal dispatch
on a $T\opt$ is rejected, the programmer has exactly two ways forward, and the
checker enforces the choice: \emph{resolve} the failure now with \verb|recovered|,
or \emph{defer} it with the fragile dispatch \verb|?.|. The rule for \verb|x?.y|
is to look past a possible $\bot$, dispatch \verb|.y| on what is inside (fully
shape-checked), and re-wrap the answer as maybe-$\bot$ --- so fragility rides
along the chain, and \verb|x?.y?.z?.a| comes out as $A\opt$, a value that is still
maybe-$\bot$ until some \verb|recovered| ends it. This is ordinary optional
chaining (Rust's \verb|?|, TypeScript's \verb|?.|); it turns ``I carried the
$\bot$ this far'' into something the compiler tracks and, at the first use as a
definite value, forces you to settle. It does not switch off checking:
\verb|x?.y?.z?.bogus| is still rejected when that shape has no \verb|bogus|. And
\verb|recovered| is checked too --- \verb|(recovered x alt).y| is accepted only
when the fallback \verb|alt| can also answer \verb|.y|, so the two arms really are
interchangeable.

\paragraph{Where $\bot$ comes from.} The checker never \emph{guesses} fragility
from how code is written. $T\opt$ enters only from declared \emph{sources} --- an
atom whose stated result is $T\opt$ (a lookup that can miss, a slice that can run
off the end), the \verb|?.| operator, and the rule just below --- and then
propagates by ordinary flow. An object is fragile in a given place exactly when
its inferred type there is $T\opt$, traced back to one of those sources. (A
corollary worth stating: today's parser output predates the $\bot$/\verb|?.|
syntax, so the real objects show almost no fragility yet --- the machinery is
exercised by hand-built cases and by the rule below.)

\paragraph{Incompleteness as fragility.} EO adds one more source of $\bot$, and a
tidy one. Some voids are optional recovery callbacks --- \verb|cant-convert| on a
numeric conversion, \verb|cant-connect| on a socket --- which a confident
programmer may leave unset. But an object with any attribute unset is
\emph{incomplete}, and an incomplete object is fragile: using it must go through
\verb|?.| or be recovered. This needs no new machinery --- a formation that
escapes as a value with voids still unfilled is given type $T\opt$; application is
allowed to see through the $\opt$ (so \verb|foo 1 2| still completes a
two-argument \verb|foo| to a definite value), while any plain dispatch on it meets
the discipline above. Run over the whole runtime with the rule on, the checker
flags exactly five objects: \verb|number| and \verb|i64| (a conversion used with
\verb|cant-convert| unset), \verb|fs/file| (\verb|cant-check|), \verb|nk/socket|
(\verb|cant-receive|), and \verb|fs/path| (an unapplied \verb|uri|). Four are
precisely the recovery-callback situations the rule is meant to catch. The fifth,
reading \verb|posix.separator|, is \emph{not} a false positive. The rule is
deliberately \emph{object-level} --- an object with any attribute unset is
fragile, full stop --- so dispatching \verb|.separator| on a \verb|posix| whose
\verb|uri| was never supplied is correctly rejected, even though \verb|separator|
reads only a constant. Reaching into a half-built object is exactly the smell the
rule exists to forbid; it is not object-oriented, and the fix belongs in the EO
code (give \verb|separator| a home that does not demand a \verb|uri|), not in a
weaker, attribute-level rule. This is the same kind of latent-smell surfacing as
the \verb|while| and \verb|switch| results below --- a type the author never
intended, pointing at a design that should change.

\section{The atoms, annotated}\label{sec:atoms}

The built-in objects are the \emph{starting facts}: the checker cannot look
inside them (they are implemented in Java), so we state their shapes by hand,
and everything else is worked out from these. Writing these down is step one of
any real effort --- if a shape cannot be stated here, the grammar of Section~\ref{sec:grammar}
is not yet good enough. The key ones:

\begin{center}\small
\renewcommand{\arraystretch}{1.15}
\begin{tabular}{@{}lll@{}}
\toprule
\textbf{EO} & \textbf{Type} & \textbf{In plain words} \\
\midrule
\verb|number.plus|      & $\ty{number}\arr\ty{number}$                                   & a number in, a number out \\
\verb|number.gt|        & $\ty{number}\arr\ty{bool}$                                     & compare; give \verb|true|/\verb|false| \\
\verb|bool.if|          & for any $A,B$: $A\arr B\arr (A\mid B)$                          & pick one of two branches \\
\verb|recovered|        & for any $A,B$: $A\opt\arr B\arr (A\mid B)$                      & the value, unless $\bot$, then the fallback \\
\verb|seq|              & for any $A$: $(\text{tuple ending in } A)\arr (A\mid\ty{bool})$   & run every step, give the last \\
\verb|tuple.head|       & for a $\ty{tuple}\,A$: $A$                                      & the last element (EO tuples grow at the head) \\
\verb|tuple.with|       & for a $\ty{tuple}\,A$: $B\arr \ty{tuple}\,(A\mid B)$       & append one element \\
\verb|map.found(k).get| & for any $B$: $B\arr (V\mid B)$                                  & the stored value, or your fallback \\
\bottomrule
\end{tabular}
\end{center}

\noindent
Two things to notice. \verb|bool.if| shows why we need ``for any $A,B$'':
the branches may have different types, and the answer is ``one or the other.''
\verb|seq|'s answer depends on \emph{which} element is last, but that is fixed
the moment the tuple is written, so no program needs to run --- the tool just
reads off the last position's type.

\section{The checker, over the \texttt{@}-core}\label{sec:checker}

The core of EO has only four moves: \textbf{make} an object (a formation),
\textbf{apply} arguments to it, \textbf{dispatch} an attribute (a call
\verb|e.m|), and \textbf{decorate} via \texttt{@}. A checker only has to handle
these four, plus looking up names. Here it is as plain steps; \verb|typeof|
returns the shape of an expression, and \verb|usable-as| answers ``can an X go
where a Y is wanted?''

\begin{lstlisting}
typeof(expr, env):                       # env maps names -> known types
  NAME  x:            return env[x]
  MAKE  [p1 p2 ...] with body:
      shape = {}
      for each free attribute p:         shape[p] = a fresh unknown
      for each named child  n = e:        shape[n] = typeof(e, env + params)
      if there is an `@' child e:          shape["@"] = typeof(e, env + params)
      return shape
  APPLY  f arg  (fills f's next free attribute p):
      ft = typeof(f, env);  at = typeof(arg, env)
      require  usable-as(at, ft.expected(p))          # else ERROR
      return  ft  with p replaced by at
  DISPATCH  e.m:
      t = typeof(e, env)
      if t is an option (T?):  ERROR "`m' used on a value that can be bottom;
                                      recover it first"
      look up m in t;  if missing, follow t["@"] and look again (repeat)
      if still missing:        ERROR "`m' is not an attribute of this object"
      return the found attribute's type

usable-as(X, Y):                         # "an X can go where a Y is wanted"
  if Y is an unknown:   set Y := X;  ok
  if X is an option and Y is not:   NOT ok           # a bottom might leak
  if both are object shapes:   ok  iff  every attribute Y requires is present
        in X (directly or through @), and usable-as holds on each of them
\end{lstlisting}

\noindent\textbf{A worked run.} Take a correct program:
\begin{lstlisting}
[] > my-num
  5 > @
my-num.plus 1
\end{lstlisting}
\begin{enumerate}\setlength{\itemsep}{1pt}
\item \verb|typeof(5)| is \verb|number| (a number literal).
\item \verb|typeof(my-num)|: a formation with an \verb|@| child, so its shape is
  $\{\,\ty{@}:\ty{number}\,\}$.
\item \verb|my-num.plus| is a \textsc{dispatch}: \verb|my-num| has no
  \verb|plus| of its own, so follow \verb|@| to \verb|number|, which has
  \verb|plus| $:\ty{number}\arr\ty{number}$. Found.
\item \textsc{apply} to \verb|1|: \verb|1| is a \verb|number|, and \verb|plus|
  wants a \verb|number| --- \verb|usable-as| holds. Result: \verb|number|.
\end{enumerate}
Accepted; the whole thing is a \verb|number|. Now two programs it rejects:
\begin{lstlisting}
[] > my-num              (m.get "k").plus 1     # m.get can fail, so its type
  "hi" > @                                      # is `number?' (number or bottom)
my-num.plus 1
\end{lstlisting}
On the left, \verb|typeof(my-num)| is $\{\ty{@}:\ty{string}\}$; \textsc{dispatch}
of \verb|plus| finds none on \verb|my-num|, follows \verb|@| to \verb|string|,
finds none there either --- \emph{error:} ``\texttt{.plus} is not an attribute of
\texttt{my-num} (it decorates a string).'' On the right, the receiver has type
\verb|number?|, so \textsc{dispatch} stops immediately --- \emph{error:}
``\texttt{.plus} used on a value that can be $\bot$ (from \texttt{m.get});
recover it first, e.g.\ \texttt{recovered (m.get "k") 0}.'' Recovering turns the
\verb|number?| back into a \verb|number| (that is \verb|recovered|'s type), and
the fixed program is accepted. Both mistakes are caught without running
anything.

\section{The spine: solving, and surviving recursion}\label{sec:spine}

The steps in Section~\ref{sec:checker} decide each fit \emph{on the spot}. Two
things break that --- the top two gaps from typing real objects by hand --- and
both are fixed by one idea used twice: \textbf{assume, then reconcile}.

\paragraph{Collect, then solve.} Split typing into two passes. \emph{Collect:}
walk the code once, give every unknown a fresh name, and instead of deciding
fits, just \emph{write down} each requirement as a note --- ``this unknown must
be usable as a \verb|number|,'' ``this one must have a \verb|.head|,'' ``this
value flows into that slot.'' \emph{Solve:} work through the notes and pin each
variable down. Why two passes? A variable often makes sense only once you have
seen all its uses: if \verb|x| is used once where a number is wanted and once
where a string is wanted, no single on-the-spot decision is right, but the two
notes together say \verb|x| is $\ty{number}\mid\ty{string}$. Deciding early
cannot produce that union; collecting can.

Solving needs almost no machinery. Give each variable two lists:
\textbf{flows-in} (things put into it) and \textbf{flows-out} (slots it is used
in). To process a note ``$X$ usable-as $Y$'': if $X$ is a variable, add $Y$ to
its flows-out; if $Y$ is a variable, add $X$ to its flows-in; either way,
re-check that everything already on the other list still fits (which may make
smaller notes); if both are object shapes, break it into one note per attribute
$Y$ asks for; and if $X$ can be $\bot$ while $Y$ cannot, that is the ``recover
it first'' error. A variable's final type is \emph{everything that flowed into
it, or-ed together}. So unions and ``for any'' are never invented by hand ---
they fall out of the bookkeeping.\footnote{This two-list flows-in/flows-out
bookkeeping is essentially Simple-sub, the readable form of algebraic subtyping
from the Section~\ref{sec:fail} footnote --- about 500 lines including a parser.}

\paragraph{Tie the knot.} Recursion --- an object that mentions itself --- made
the Section~\ref{sec:checker} walk loop forever (\verb|seq|, \verb|map|,
\verb|switch|, \verb|while|, and \verb|number| all do it). Same idea, on the
object: before looking \emph{inside}, give it a fresh variable \verb|t| and
remember ``this object $=$ \verb|t|.'' Now type the body; every self-reference
finds \verb|t| --- a name, not another trip inside --- so the walk stops. The
body works out to a shape $S$; close the loop by recording \verb|t|~$=S$. Since
$S$ may contain \verb|t|, that is a self-referring shape --- the ``loop
$A.\dots$'' from Section~\ref{sec:grammar}. One last loop hides in the solver:
two self-referring types can bounce a note back and forth forever, so the solver
keeps a \textbf{seen} set and stops when a note comes back. Assume-then-reconcile,
one level down.

Here is Section~\ref{sec:checker}'s checker with the spine in place ---
\verb|infer| returns a type (with variables) and piles up notes; \verb|solve|
settles them:
\begin{lstlisting}
infer(expr, env):
  NAME x:   return env[x]
  MAKE  [free...] with children and maybe an `@':
      t = fresh variable
      env2 = env + (this object -> t)       # register BEFORE descending
      for each free attribute p:  shape[p] = fresh variable
      for each child  n = e:      shape[n] = infer(e, env2 + params)
      if there is an `@' child e:  shape["@"] = infer(e, env2 + params)
      note( t = shape )                     # tie the knot
      return t
  APPLY  f arg:
      r = fresh variable
      note( infer(f,env) usable-as (takes infer(arg,env), gives r) )
      return r
  DISPATCH  e.m:
      r = fresh variable
      note( infer(e,env) usable-as { m : r, ... } )
      return r

solve(note "X usable-as Y", seen):
  if (X,Y) in seen: return                  # stop cycles (recursive types)
  seen.add (X,Y)
  if X is a variable:  X.flowsOut += Y;  for L in X.flowsIn:  solve(L usable-as Y)
  elif Y is a variable: Y.flowsIn += X;  for U in Y.flowsOut: solve(X usable-as U)
  elif both are shapes: for each label m in Y:  solve(X[m] usable-as Y[m])
  elif X can be bottom and Y cannot:  ERROR "recover the bottom first"
# afterwards: each variable's type = the union of its flows-in
\end{lstlisting}

\paragraph{A recursive run that used to loop.} Take a small self-referring
object:
\begin{lstlisting}
[tup] > last
  if. > @
    tup.tail.length.eq 0
    tup.head
    last tup.tail
\end{lstlisting}
\begin{enumerate}\setlength{\itemsep}{1pt}
\item Start on \verb|last|: fresh \verb|t|, remember \verb|last|~$=$~\verb|t|.
\item Walk the body. \verb|tup| is a fresh \verb|u|. \verb|tup.head| notes
  ``\verb|u| has \verb|head| of type \verb|w|''; \verb|tup.tail| notes ``\verb|u|
  has \verb|tail| of type \verb|v|.''
\item \verb|last tup.tail| is the self-call. \verb|last| is \verb|t| --- a
  \emph{variable} (from step 1), \emph{not} another descent --- so the walk
  \textbf{stops here}, noting ``\verb|t| takes \verb|v|, gives \verb|r|.''
\item \verb|if. cond A B| gives \verb|A|~$\mid$~\verb|B|, here
  \verb|w|~$\mid$~\verb|r|. Close the knot: \verb|t|~$=$~(\verb|u|~$\to$~(\verb|w|~$\mid$~\verb|r|)),
  self-referring through \verb|r|.
\item Solving settles the notes: \verb|u| is a tuple, and \verb|w|, \verb|r| are
  its element type, giving \verb|last| the type ``for any $A$:
  $\ty{tuple}\,A \to A$''. No loop, no annotations. (That ``\verb|tail| of a
  tuple is the same tuple'' step needs the tuple rules --- the next piece to add.)
\end{enumerate}
The spine closes gaps~1 and~2. The remaining Tier-1 items --- tuple rules,
applying the ``for any'' atom signatures, and pushing types into callbacks ---
build directly on this collect-and-solve frame.

\section{Completing the core: tuples, generic atoms, callbacks}\label{sec:core-rest}

With the spine (Section~\ref{sec:spine}) in place, the other three Tier-1 gaps
are each a small addition to the same collect-and-solve frame.

\paragraph{Tuples.} Give a tuple two possible shapes: a \emph{fixed} one,
$\langle\tau_1,\dots,\tau_n\rangle$, when the length is known (as in
\verb|* a b c|), and a \emph{loose} one, $\ty{tuple}\,A$ (``every element is
usable as $A$''), when it is not (as when a tuple is grown in a loop). The
accessors follow --- recalling that EO tuples grow at the head, so \verb|head|
is the \emph{last} element:
\begin{itemize}\setlength{\itemsep}{0pt}\setlength{\parskip}{0pt}
\item \verb|* a b c| $:\;\langle\tau_a,\tau_b,\tau_c\rangle$; \quad
      \verb|t.length| $:\;\ty{number}$.
\item \verb|t.head| $:\;\tau_n$ (fixed) or $A$ (loose); \quad
      \verb|t.tail| $:\;\langle\tau_1,\dots,\tau_{n-1}\rangle$ or $\ty{tuple}\,A$.
\item \verb|t.with x| $:\;\langle\tau_1,\dots,\tau_n,\tau_x\rangle$ or
      $\ty{tuple}\,(A\mid\tau_x)$.
\end{itemize}
The solver already breaks a shape into one note per part, and a tuple is just
another shape: ``$X$ usable-as $\langle\dots\rangle$'' becomes one note per
position (fixed) or one note on the shared element (loose). The fixed form is
what lets \verb|seq| and \verb|switch| pin their result to the \emph{last}
element without running anything; a loose tuple falls back to the union of its
elements.

\paragraph{Generic atoms: instantiate on use, generalize on definition.} The
Section~\ref{sec:atoms} signatures are \emph{schemes} --- types with ``for any''
holes --- and two moves connect them to the checker. \textbf{Instantiate:} at
each use, make \emph{fresh} copies of the ``for any'' variables, then note the
arguments against them. \verb|recovered 42 7| makes a fresh \verb|A| and
\verb|B|, notes ``\verb|42| flows into \verb|A|?'' (so \verb|A| becomes
\verb|number|) and ``\verb|7| flows into \verb|B|'' (\verb|B| becomes
\verb|number|), and the result \verb|A| $\mid$ \verb|B| solves to \verb|number|
--- an argument-dependent result computed by the solver, no evaluation, and the
fresh copies stop two calls at different types from tangling. \textbf{Generalize:}
when a user object's inferred type still has unknowns that nothing outside
constrains, freeze them as ``for any,'' turning a one-off type into a reusable
scheme so \verb|reduced| works at every element type. These are the two standard
moves of Hindley--Milner, with each variable's flows-in/flows-out bounds riding
along.

\paragraph{Making it sound: levels.} One question decides whether generalization
helps or hurts: \emph{which} unknowns may be frozen as ``for any''? Freeze one
that is really shared with the surrounding object and you duplicate it wrongly;
freeze too few and a reused object collapses all its call sites into one --- the
path object \verb|fs/path| uses its helper \verb|posix| at a dozen sites, and
without generalization their arguments merge, so a parameter that is a string at
every real call is inferred as an impossible \verb|string|~\&~\verb|bytes| and a
correct object is rejected. The standard cure is to stamp every unknown with a
\textbf{level} --- the nesting depth of the object being defined --- and, at a
use, to instantiate exactly the variables \emph{deeper} than that use while
sharing the rest; a companion step, \emph{extrusion}, lowers the level of any
variable that escapes to a shallower scope so that nothing still-shared is ever
frozen. Levels are what make ``instantiate on use, generalize on definition''
sound rather than merely plausible, and they are the part the prototype in the
next section spent the most care getting right.

\paragraph{Callbacks: push the expected type in.} An atom that takes a function
says so in its signature --- \verb|malloc.for|'s scope is \verb|block -> R|,
\verb|while|'s condition is \verb|number -> bool|. When the argument is a
formation \verb|[m] ...|, do not type it blind: \emph{check} it against that
expected function type, which seeds the parameter (\verb|m| is a \verb|block|)
before typing the body --- so \verb|m.put| and \verb|m.read| resolve, and any
mistake is reported right at the callback. Having the checker both \emph{infer}
a type and \emph{check against} an expected one is the only real addition here.
One rule makes it work: to fit one function where another is wanted, the results
go the same way but the \emph{arguments go the opposite way} --- a
\verb|block -> R| is wanted, so the callback's parameter must be willing to
\emph{accept} a \verb|block|. That backwards flow is what carries the
\verb|block| into \verb|m|, and it is the one genuinely counter-intuitive rule
in the whole system.\footnote{Both \emph{inferring} and \emph{checking against}
an expected type is \emph{bidirectional typing}: a friendly survey is
J.~Dunfield \& N.~Krishnaswami, \emph{Bidirectional typing}, ACM Computing
Surveys, 2021. ``Arguments go the opposite way'' is standard function subtyping
--- contravariant in the argument, covariant in the result.}

With the spine and these three additions, the checker can walk the whole
standard library --- which is what the next section shows.

\section{It scales to real code: a working checker}

These rules are not only on paper. A prototype of about 900 lines implements
exactly this --- the grammar, the annotated atoms, the collect-then-solve spine,
and the level-based generalization above --- and reads EO's \emph{real} parser
output (the XMIR under \texttt{eo-runtime/target/eo/1-parse}). Run over the whole
runtime it types \textbf{all 117 top-level objects}, with no rejections, in under
a second. And --- the actual test of a type system --- it still \emph{catches}
mistakes: a string sent \verb|.plus|, a value used at once as a function and as a
string, a dispatch on an attribute that is not there are all reported, and every
deliberately-wrong example in Sections~\ref{sec:fail}--\ref{sec:core-rest}
rejects as designed. Passing everything is trivial for a checker that says yes to
everything; this one says no in the right places.

The standard library comes out as the rules predict. The list combinators
(\texttt{tuple/mapped.eo}, \texttt{tuple/reduced.eo}) are generic: \verb|mapped|
has type ``for any $A,B$: $\ty{tuple}\,A \arr (A\arr B) \arr \ty{tuple}\,B$'' and
\verb|reduced| is the usual fold. Their callbacks carry no annotations --- the
tool reads each callback's type from its body --- so
\verb|mapped (* 1 2 3) (x.times 2 > [x])| is a tuple of numbers, while a callback
returning a \verb|bool| makes a later \verb|.plus| on an element a caught error.
A \verb|map| never returns a bare value: \verb|found(key).get| \emph{demands} a
fallback, so a lookup has type ``value \emph{or} fallback'' --- absence lives in
the type. And \verb|switch| returns \verb|true| when no case matches, so its
inferred type carries a stray \verb|bool| the author never intended: the type
quietly points at a real bug no test triggers as long as some case matches.

\paragraph{What took the care.} Almost all of it was the level machinery, and
four subtleties were decisive --- each a place the obvious implementation is
wrong. The self-reference that ties a recursive knot must live \emph{inside} the
object's own scope, not outside it, or generalization silently de-generalizes the
parameters and the checker stops catching mistakes. The built-in atoms must sit
at the outermost level, or every use of one drags the solver into needless
extrusion and it crawls. The parent link $\rho$ must be left out when measuring
an object's level, or a deep parent keeps a record ``too deep'' to extrude and
solving never stops. And a large shared type must be printed by naming each
variable once, or \verb|fs/path|'s type --- a wide, cyclic shape --- prints
exponentially. None of these touches the design of
Sections~\ref{sec:grammar}--\ref{sec:core-rest}; all are the ordinary hazards of
a level-based solver, now charted.

\section{Is this doable, and your ``just compile it'' idea}

Yes, and it is not exotic: a type based on shapes, worked out automatically,
with ``usable-as'' and ``$A$ or $B$,'' is a well-charted design.\footnote{The
combination --- structural shapes, ``usable-as,'' unions, and full inference ---
is the subject of the algebraic-subtyping work cited in Section~\ref{sec:fail}; a gentle
textbook is B.~Pierce, \emph{Types and Programming Languages}, MIT Press 2002.}
It belongs in a compile-time step over XMIR (EO's intermediate form), never in
the runtime.

Your idea of translating EO into a strictly-typed language (say Haskell) and
letting its compiler judge is a real technique, and half of it fits EO
perfectly: $\bot$ maps to Haskell's \texttt{Maybe} and \verb|recovered| to
\texttt{fromMaybe}. The catch is the other half. Haskell decides types by
\emph{name} and has no built-in ``usable-as,'' which is precisely EO's
decoration --- so you would rebuild the hard part inside Haskell's most awkward
corner, the errors would come out talking about the generated Haskell rather
than the user's EO, and ``if it compiles, it is safe'' is only true if the
translation is faithful, which you can only know by writing down the rules in
Sections~\ref{sec:grammar}--\ref{sec:core-rest} first. If you do want to borrow a compiler, pick a
\emph{structurally}-typed one: TypeScript already decides types by shape and has
$A\mid B$ and ``value or nothing'' built in --- almost the whole EO picture ---
and OCaml's object types are the same idea with more rigour. Best use: a quick
prototype to test these rules, not the finished checker. Either way the rules
come first.

So the order is: (1)~this grammar; (2)~these atom shapes; (3)~this checker,
first on the four-move core, then widened, dog-fooded on the standard library.
The parts that will take real care are self-referring shapes (the ``loop''
case), keeping the tool's guesses predictable, and
error messages that speak EO. None is unexplored; all are documented. EO is
ready to start --- the hard idea, that decoration forces us to talk about shape,
is already behind us.\footnote{EO and its calculus: Y.~Bugayenko, \emph{EOLANG
and $\varphi$-calculus}, arXiv:2111.13384, 2021; N.~Kudasov \& V.~Sim,
\emph{Formalizing $\varphi$-calculus}, arXiv:2204.07454, 2022.}

\end{document}
